\subsection{序域}

定义 2.——交换域配备一个全序，若该序与其环结构相容，则称为序域。

我们仅考虑域上的全序关系，因为其他的序关系非常“病态”（VI，第 38 页，习题 6）。

例子。—— 1）有理数域 \(\mathbf{Q}\) 是一个序域。

\begin{enumerate}
\def\labelenumi{\arabic{enumi})}
\setcounter{enumi}{1}
\item
  序域的子域，赋予诱导序，也是一个序域。
\item
  * 实数域是一个序域。 *
\end{enumerate}

设 \(\mathbf{K}\) 是一个序域。对所有 \(x\in \mathbf{K}\) 我们令

\[
\begin{array}{l l} \operatorname{sgn} (x) = 1 & \text {if} \quad x > 0, \\ \operatorname{sgn} (x) = - 1 & \text {if} \quad x <   0, \\ \operatorname{sgn} (x) = 0 & \text {if} \quad x = 0. \end{array}
\]

那么我们有 \(\operatorname{sgn}(xy) = \operatorname{sgn}(x)\operatorname{sgn}(y)\) ；我们称 \(\operatorname{sgn}(x)\) 为 \(x\) 的符号。映射 \(x \mapsto \operatorname{sgn}(x)\) 从 \(\mathbf{K}^*\) 到乘法群 \(\{-1, +1\}\) 是一个满射同态，其核，即 \(\mathbf{K}\) 的严格正元素所成之集，是 \(\mathbf{K}^*\) 的一个指数为 2 的子群。

反之，若 K 是一个交换域，且 \(s: K^{*} \to \{-1, +1\}\) 是一个核在加法下封闭的满同态，则 s 是某个唯一的序域结构的符号映射，其中严格正元素之集就是 s 的核。

对所有 \(x\) 与 \(y\) 属于 \(\mathbf{K}\) ，我们有 \(x = \operatorname{sgn}(x)|x|\) 且 \(|xy| = |x||y|\) 。

另一方面，每个序域的特征都为零（命题 1）。

命题 2. --- 设 A 是一个全序整环，K 是它的分式域。则 K 上存在唯一一个序，它在 A 上限制为给定的序，并使 K 成为一个序域。

每个 \(x \in K\) 都可表为 \(x = ab^{-1}\) ，其中 a、b 属于 A 且 \(b \neq 0\) 。若 x 为正，则 a 与 b 同号，反之亦然。由此看出，若 K 上存在满足所述条件的序，则它是唯一的，且正元素之集 P 恰为所有 \(ab^{-1}\) 的集合，其中 a、b 是 A 中同号的元素且 \(b \neq 0\) 。剩下要证明 P 满足条件 \((\mathrm{AP}_{\mathrm{I}}), (\mathrm{AP}_{\mathrm{II}}), (\mathrm{AP}_{\mathrm{III}})\) 与 \((\mathrm{AP}_{\mathrm{IV}})\) 。这对 \((\mathrm{AP}_{\mathrm{II}})\) 与 \((\mathrm{AP}_{\mathrm{IV}})\) 是显然的。对 \((\mathrm{AP}_{\mathrm{I}})\) ，考虑 \(ab^{-1} + cd^{-1}\) ，其中可设 \(a, b, c\) 与 \(d\) 为正；这个和是 \((ad + bc)(bd)^{-1}\) ，且 \(ad + bc\) 与 \(bd\) 为正。

为证 \((\mathrm{AP}_{\mathrm{III}})\) ，考虑一个形如 \(ab^{-1} = -cd^{-1}\) 的恒等式，从而 \(ad + bc = 0\) 。若设 a 与 b 同号且 c 与 d 同号，则符号法则表明 ad 与 bc 同号；于是 ad = bc = 0，故 a = c = 0；从而 P 确实满足 \((\mathrm{AP}_{\mathrm{III}})\) 。

例。--- 由于 Z 只容许一个全序环结构（VI，第 19 页，例），域 Q 只容许一个使它成为序域的序：这就是通常的序。

